\section{Introducción}

Una planta solar fotovoltaica genera, a través de sus inversores, un volumen continuo de mediciones eléctricas (corriente y voltaje DC de entrada, corriente, voltaje y potencia AC de salida) que reflejan tanto las condiciones ambientales (irradiancia solar) como el estado operativo del equipo. Analizar estos datos con técnicas de minería de datos permite pasar de un registro pasivo de mediciones a información accionable para tres decisiones de negocio concretas: planeación del despacho eléctrico, mantenimiento predictivo de inversores y pronóstico de producción.

En este trabajo se analiza el dataset \texttt{2107\_electrical\_data.csv}, que contiene mediciones cada 5 minutos de los 24 inversores de una planta solar durante aproximadamente 6 años (noviembre de 2017 a noviembre de 2023). No se proporcionó un problema específico a resolver: el conjunto fue explorado para identificar oportunidades de análisis, resultando en tres líneas de trabajo (supervisado, no supervisado y forecasting) descritas a continuación.

\subsection*{Planteamiento del problema}

El dataset no incluye variables meteorológicas externas (irradiancia, temperatura ambiente), por lo que toda la información disponible para diagnóstico y predicción proviene únicamente de las señales eléctricas internas de la planta. Esto plantea tres preguntas concretas:

\begin{enumerate}
    \item ¿Es posible predecir la potencia AC total de la planta a partir de las señales de entrada de los inversores, con precisión suficiente para apoyar decisiones operativas? (\textbf{aprendizaje supervisado})
    \item ¿Existen inversores con un perfil de desempeño distinto al resto de la flota, o intervalos de tiempo con comportamiento anómalo que sugieran una falla o degradación? (\textbf{aprendizaje no supervisado / diagnóstico})
    \item ¿Puede anticiparse la energía que producirá la planta en las próximas semanas a partir de su historial? (\textbf{forecasting})
\end{enumerate}

\subsubsection*{Objetivo general}

Aplicar un proceso completo de ciencia de datos —comprensión, exploración, preparación, modelado y evaluación— sobre el dataset eléctrico de la planta solar, para diagnosticar el comportamiento de sus variables y generar modelos predictivos que apoyen la toma de decisiones.

\subsubsection*{Objetivos específicos}

\begin{itemize}
    \item Caracterizar la calidad, completitud y comportamiento temporal de las señales eléctricas de los 24 inversores.
    \item Preparar y justificar las transformaciones necesarias para el modelado (imputación, variables derivadas, escalamiento).
    \item Entrenar y comparar modelos de aprendizaje supervisado para predecir la potencia AC total de la planta.
    \item Aplicar técnicas no supervisadas para segmentar inversores por desempeño y detectar anomalías operativas.
    \item Pronosticar la producción energética futura de la planta mediante un modelo de series de tiempo.
    \item Discutir los resultados y sus limitaciones, y sintetizar las aportaciones del trabajo.
\end{itemize}

\subsection*{Metodología utilizada}

\begin{itemize}
    \item \textbf{Carga y exploración.} Lectura del CSV con \texttt{pandas}, revisión de tipos de dato, nulos, duplicados y rango temporal.
    \item \textbf{Preparación.} Imputación justificada de nulos, corrección de un error tipográfico en un nombre de columna, creación de variables derivadas (potencia total, variables temporales) y escalamiento (\texttt{StandardScaler}).
    \item \textbf{Aprendizaje supervisado.} Comparación de Regresión Lineal, Random Forest y un Perceptrón Multicapa (MLP) entrenado con \texttt{PyTorch} (optimizador Adam, pérdida MSE).
    \item \textbf{Aprendizaje no supervisado.} Perfilado de inversores, reducción de dimensión con PCA, agrupamiento con K-Means y detección de anomalías con Isolation Forest.
    \item \textbf{Forecasting.} Descomposición estacional y suavizado exponencial de Holt-Winters sobre la serie semanal de energía generada, contra un baseline ingenuo estacional.
    \item \textbf{Documentación.} Notebook \texttt{plan/plan\_mineria.ipynb} y este reporte en \LaTeX{}.
\end{itemize}
